\cleardoublepage 
\chapter{PENDAHULUAN}
\label{chap:pendahuluan}

\section{Latar Belakang}
Tindak kriminal pada ruang digital yang memanfaatkan \textit{artificial intelligence} (AI) untuk memanipulasi audio/suara manusia kian meningkat dan sangat memprihatinkan. AI dimanfaatkan para pelaku kejahatan untuk membuat \textit{deepfake audio} yang ditujukan kepada korban sebagai modus penipuan atau pemerasan. Kebanyakan korban teperdaya oleh hasil manipulasi tersebut, yang mengakibatkan kerugian pada banyak aspek, khususnya kerugian finansial.

Para pelaku kejahatan melakukan hal tersebut dengan tujuan untuk menutup identitas aslinya. Metode tersebut berpotensi menimbulkan fitnah digital, manipulasi alat bukti elektronik, hingga kesalahan identifikasi dalam proses penegakan hukum \autocite{Iman_2026}. Upaya penegakkan hukum untuk memberantas permasalahan ini memerlukan bantuan forensik digital. Forensik digital diperlukan dalam memperoleh bukti-bukti digital, termasuk bukti \textit{deepfake audio} yang dibuat AI. Namun, hal ini hanya dapat dilaksanakan pada laboratorium forensik digital yang memfasilitasi proses perolehan serta pendeteksian bukti digital oleh para ahli.

Saat ini, laboratorium forensik digital hanya ada di Dittipidsiber Mabes Polri dan Polda Metro Jaya. Selain itu, keterbatasan fasilitas tersebut pun diiringi dengan kurangnya ahli forensik digital, khususnya ahli yang memahami AI generatif \autocite{SIP_Law_Firm_2025}. Keterbatasan ini menimbulkan kesulitan bagi para korban yang ingin melakukan pelaporan atau pembuktian. Korban harus merekam audio sebagai bukti sah yang selanjutnya dilimpahkan kepada institusi terkait untuk diproses.

Oleh karena itu, diperlukan suatu metode yang mempermudah korban untuk mendeteksi \textit{deepfake audio} buatan AI. Para penulis menggagas platform bertajuk \textit{Audio Deepfake Anomaly Detection Harness} (ADADeH) untuk memerangi manipulasi AI dengan AI (AI vs AI) yang dilatih untuk mendeteksi \textit{deepfake audio}. Dengan demikian, pembuktian tidak perlu melalui serangkaian prosedur yang rumit.

\section{Rumusan Masalah}
Keterbatasan sarana pembuktian yang ada menimbulkan beberapa persoalan dalam menghadapi kejahatan berbasis \textit{deepfake audio}. Proses pendeteksian saat ini bergantung pada laboratorium forensik digital dan tenaga ahli yang jumlahnya terbatas, sehingga korban sulit mendapatkan indikasi awal atas rekaman yang diterimanya. Selain itu, suara hasil AI generatif kian menyerupai suara asli manusia sehingga tidak dapat dibedakan hanya dengan pendengaran. Di sisi lain, hasil deteksi yang ada umumnya belum disajikan dalam bentuk yang mudah dipahami masyarakat awam.

Apabila persoalan-persoalan ini tidak segera ditangani, korban akan tetap kesulitan melakukan pelaporan dan pembuktian, sementara pelaku dapat terus memanfaatkan \textit{deepfake audio} untuk menutupi identitasnya. Oleh karena itu, tugas proyek ini mengusulkan ADADeH, yaitu platform berbasis AI yang mendeteksi anomali pada rekaman audio untuk membedakan suara asli manusia dari \textit{deepfake audio}. 

Berdasarkan hal tersebut, rumusan masalah setidaknya mencakup hal-hal berikut.
\begin{enumerate}
\item	Bagaimana mendeteksi anomali pada rekaman audio yang membedakan suara asli manusia dari \textit{deepfake audio} hasil AI generatif?
\item	Bagaimana merancang \textit{pipeline} ADADeH, mulai dari ekstraksi fitur audio hingga klasifikasi, dan evaluasi kinerja menggunakan metrik yang sesuai?
\item	Bagaimana menyajikan hasil deteksi ADADeH secara mudah sebagai indikasi awal pembuktian, tanpa bergantung pada laboratorium forensik digital?
\end{enumerate}

\section{Tujuan}
Berdasarkan latar belakang dan rumusan masalah yang telah diuraikan, tujuan dari TPdK ini adalah mengembangkan ADADeH, yaitu \textit{platform} berbasis AI untuk mendeteksi \textit{deepfake audio} sehingga dapat memberikan indikasi awal pembuktian bagi korban tanpa bergantung pada laboratorium forensik digital. Secara rinci, tujuan tersebut dijabarkan sebagai berikut.
\begin{enumerate}
\item	Membangun model deteksi anomali yang dapat membedakan suara asli manusia dari \textit{deepfake audio} hasil AI generatif berdasarkan karakteristik akustik rekaman.
\item	Merancang dan mengimplementasikan \textit{pipeline} ADADeH yang mencakup ekstraksi fitur audio, klasifikasi, dan evaluasi kinerja menggunakan metrik yang sesuai.
\item	Menyajikan hasil deteksi dalam bentuk yang mudah dipahami orang awam sebagai indikasi awal pembuktian.
\end{enumerate}

Tugas proyek ini dinyatakan berhasil apabila ketiga tujuan tersebut tercapai. Kriteria keberhasilan yang digunakan adalah sebagai berikut.
\begin{enumerate}
\item	Model deteksi mencapai akurasi dan \textit{F1-score} minimal 90\% serta \textit{Equal Error Rate} (EER) maksimal 10\% pada data uji yang tidak digunakan dalam proses pelatihan.
\item	\textit{Pipeline} ADADeH dapat dijalankan secara utuh, mulai dari masukan rekaman audio hingga keluaran hasil klasifikasi, tanpa adanya intervensi manual.
\item	Hasil deteksi ditampilkan dalam bentuk yang dapat dipahami oleh pengguna awam, yaitu label klasifikasi beserta tingkat keyakinan model.
\end{enumerate}

% --- Lingkup dan Batasan ---
\section{Lingkup dan Batasan}
Lingkup dan batasan pengembangan ADADeH adalah sebagai berikut.
\begin{enumerate}
\item ADADeH membantu korban kejahatan berbasis \textit{deepfake audio} memperoleh indikasi awal atas keaslian suatu rekaman. Sistem mencakup pemrosesan rekaman audio, ekstraksi fitur, klasifikasi suara asli dan suara hasil AI generatif, serta penyajian hasil deteksi kepada pengguna.
\item Pelatihan dan pengujian model menggunakan satu \textit{dataset} publik Kaggle berisi rekaman berbahasa Inggris. Kinerja sistem pada teknik \textit{deepfake} lain di luar \textit{dataset} dan pada rekaman berbahasa Indonesia belum dapat dipastikan.
\item Sistem hanya melakukan klasifikasi biner untuk membedakan suara asli dan \textit{deepfake audio}, tanpa mengidentifikasi pelaku maupun teknik pembuatan \textit{deepfake} yang digunakan.
\item Masukan sistem berupa rekaman audio dalam format digital dengan kualitas rekaman yang memadai. Rekaman dengan \textit{noise} yang sangat tinggi atau durasi yang terlalu singkat dapat menurunkan akurasi deteksi.
\item Hasil deteksi tidak menggantikan pemeriksaan forensik digital oleh ahli di laboratorium forensik digital.
\item Pengembangan dibatasi hingga integrasi dan pengujian purwarupa. Penerapan dalam skala luas dan pemeliharaan sistem belum menjadi bagian dari tugas proyek ini.
\end{enumerate}

% --- Metodologi Pengembangan ---
\section{Metodologi Pengembangan}
Pengembangan ADADeH menggunakan model \textit{waterfall} dari analisis kebutuhan hingga integrasi dan pengujian. Keluaran setiap tahap menjadi dasar tahap berikutnya. Hasil akhirnya adalah purwarupa yang dapat dijalankan dan didemonstrasikan.

Selain itu, pendekatan \href{https://public.dhe.ibm.com/software/analytics/spss/documentation/modeler/14.2/es/CRISP-DM.pdf}{\textit{Cross-Industry Standard Process for Data Mining} (CRISP-DM)} digunakan untuk memandu penyiapan data audio, pemodelan, dan evaluasi model sebelum diintegrasikan ke dalam purwarupa. Jika hasil evaluasi belum memenuhi target, data atau model disempurnakan sebelum integrasi.

\begin{enumerate}
\item \textbf{Analisis kebutuhan.} Masalah \textit{deepfake audio} dan kebutuhan pengguna ditelaah melalui latar belakang serta literatur. Kebutuhan fungsional dan nonfungsional dirumuskan, lalu sedikitnya tiga alternatif solusi dibandingkan dengan matriks keputusan.
\item \textbf{Perancangan.} Arsitektur purwarupa dan aliran data dirancang dari masukan audio, prapemrosesan, ekstraksi fitur, dan klasifikasi hingga keluaran label serta tingkat keyakinan dalam teks. Perancangan juga mencakup pemisahan data latih dan uji serta antarmuka bagi pengguna awam.
\item \textbf{Implementasi.} Rekaman berbahasa Inggris dari \href{https://www.kaggle.com/datasets/adarshsingh0903/audio-deepfake-detection-dataset}{\textit{Audio Deepfake Detection Dataset}} di Kaggle disiapkan. Komponen pemrosesan audio, model klasifikasi suara asli dan \textit{deepfake audio}, serta antarmuka purwarupa dibangun sesuai rancangan.
\item \textbf{Integrasi dan pengujian.} Komponen digabungkan menjadi alur otomatis dari rekaman hingga hasil klasifikasi. Pada data uji yang tidak digunakan untuk pelatihan, matriks konfusi digunakan untuk menghitung akurasi dan \textit{F1-score}. Model ditargetkan mencapai kedua metrik tersebut minimal 90\% serta \textit{Equal Error Rate} (EER) maksimal 10\%. Evaluasi pengguna menilai keterpahaman label dan tingkat keyakinan.
\end{enumerate}

% --- Sistematika Penulisan ---
\section{Sistematika Penulisan}
Laporan tugas proyek ADADeH disusun dalam enam bab sebagai berikut.
\begin{enumerate}
\item \textbf{Bab I Pendahuluan.} Bab ini memuat latar belakang, rumusan masalah, tujuan, lingkup dan batasan, metodologi pengembangan, serta sistematika penulisan.
\item \textbf{Bab II Dasar Teori.} Bab ini membahas konsep sistem multimedia, karakteristik audio dan \textit{deepfake audio}, dasar pemrosesan dan deteksi audio, serta kajian solusi atau penelitian terkait.
\item \textbf{Bab III Analisis Kebutuhan.} Bab ini menjelaskan kebutuhan pengguna dan sistem, sedikitnya tiga alternatif solusi, serta dasar pemilihan solusi yang dikembangkan.
\item \textbf{Bab IV Desain dan Implementasi.} Bab ini menguraikan arsitektur ADADeH, aliran pengolahan audio, rancangan antarmuka, dan implementasi purwarupa.
\item \textbf{Bab V Pengujian dan Diskusi.} Bab ini menyajikan pengujian sistem dan model, evaluasi pengguna, serta pembahasan hasil dan keterbatasannya.
\item \textbf{Bab VI Kesimpulan dan Saran.} Bab ini menyimpulkan pencapaian tujuan ADADeH dan menyampaikan saran pengembangan.
\end{enumerate}
